\documentclass[10pt, a4paper]{article}

\usepackage[margin=1.5cm]{geometry}
\usepackage[utf8]{inputenc}
\usepackage[spanish,es-noshorthands]{babel}
\usepackage{amssymb, amsmath, amsbsy}
\usepackage{verbatim}
\usepackage{mathpazo}
\usepackage{tasks}
\usepackage{enumerate}
\usepackage[pdftex]{hyperref}
\hypersetup{pdftitle={Examen de Introducción a la Computación / Computación I },pdfauthor={Luis Carrillo Gutiérrez}}

\fontfamily{cmss}\selectfont
\renewcommand{\familydefault}{\sfdefault}
\pagestyle{empty}

\begin{document}
\begin{center}
\subsection*{Examen de la $\text{2}^{\text{da}}$ Unidad de Aprendizaje / Contabilidad} 
\subsubsection*{Computación I – UPAC / 2011 - II} 
\end{center}

\noindent Nombres / Apellidos : \rule{10cm}{0.4pt} \\ 
Fecha : \rule{2cm}{0.4pt}  / \rule{2cm}{0.4pt} / \rule{2cm}{0.4pt} \hspace{3.4cm} Duración: {\small 35 minutos} \\

\begin{enumerate}
\item{El conjunto de nodos capaces de {\tt NO} comunicarse entre sí, con servicios especializados que conmutan datos entre los participantes es un concepto de :
\begin{tasks}(2)
\task{Una red informática}
\task{Un diagrama de comunicación}
\task{Un protocolo no entendible}
\task{Ninguna de las anteriores}
\end{tasks}
}
\item{El {\tt B2C} se refiere a:
\begin{tasks}
\task{La relación entre consumidor y la empresa}
\task{La relación docente/estudiante}
\task{La relación entre jefe y empleado}
\task{Todas las anteriores}
\end{tasks}
}
\item{¿Cuál proposición es FALSA?
\begin{tasks}
\task{Los recursos conceptuales administran los recursos físicos}
\task{La información relevante es cuando es: precisa, oportuna, completa y valorizable}
\task{Un dato es mejor cuando incumplen con la REGLA de GIGO}
\task{Los programadores son los mejores a cumplir el cargo de un CIO}
\task{El ERP es una evolución coherente de lo que es un MRP II}
\end{tasks}
}
\item{Un joven nacido en el año 1993 o posteriores es considerado:
\begin{tasks}(2)
\task{Un inmaduro}
\task{Un inmigrante digital}
\task{Alguien de una generación pérdida}
\task{Un nativo digital}
\end{tasks}
}
\item{Las impresoras, cámaras de vigilancia o terminales tontos dentro de una {\tt red informática} son llamados :
\begin{tasks}(2)
\task{Nodos comunicacionales}
\task{Recursos}
\task{Elementos digitales de la red}
\task{Terminales conductoras de señal}
\end{tasks}
}
\item{¿En cuál de los siguientes ejemplos (URI) se ha implementado un servidor llamado {\tt academico}, donde se encuentra alojado un sitio que maneja \emph{hipertexto} y es de tipo \emph{educativo} alojado en Perú.
\begin{tasks}
\task{hipertext://academico.peru.educativo.com}
\task{academico://peru.educativo.edu}
\task{http://academico.uaustral.edu.pe}
\task{peru://educativo.academico}
\end{tasks}
}
\item{Un lugar donde los elementos TIC involucrados promueven la convivencia de un grupo de personas que lo habita es denominado : 
\begin{tasks}(2)
\task{Sociedad de Información}
\task{Aldea global}
\task{Sociedad del Conocimiento}
\task{Todas las anteriores}
\end{tasks}
}
\item{¿Cuál de las siguientes proposiciones es VERDADERA?
\begin{tasks}
\task{El IPv4 es mucho más conveniente a las necesidades actuales que el IPv6}
\task{Las universidades han adquirido muchos de los dominios .xxx}
\task{Un nuevo dominio .com es mucho más costoso que un dominio .info}
\task{El Perú es una de las sociedades del conocimiento más importante de Latinoamérica}
\end{tasks}
}
\item{El manejo correcto del contenido y de los meta datos relacionados a los sitios web y sus diferentes páginas web para que obtengan un mejor rankeo dentro de los motores de búsqueda se denomina :
\begin{tasks}
\task{Posicionamiento en redes sociales}
\task{Search Engine Optimization (SEO)}
\task{Construcción de archivos HTML}
\task{Marketing Offline} 
\end{tasks}
}
\item{El {\tt SaaS} es el acrónimo para el concepto de :
\begin{tasks}(2)
\task{Internet Protocol versión 4.0}
\task{Software como servicio}
\task{Una red social popular}
\task{Un tipo de SmartMob}
\end{tasks}
}
\item{Un {\tt e-Business}, transforma procesos claves del negocio, mediante el uso de :
\begin{tasks}(2)
\task{TIC}
\task{Redes}
\task{Personas}
\task{La internet}
\end{tasks}
}
\item{''Asociar diversos sistemas computacionales, para un trabajo conjunto y generar información`` es concepto de :
\begin{tasks}(2)
\task{Escalabilidad}
\task{Disponibilidad}
\task{Unicuomaniabilidad}
\task{Interoperabilidad}
\end{tasks}
}
\item{ Definir/Desarrollar los siguientes conceptos :\quad TPS/AIS,\quad Cadena de Valor,\quad Interoperabilidad,\quad El dominio .mobi,\quad Motor de búsqueda,\quad Google Page Rank,\quad CIO,\quad ERP.
}
\end{enumerate}
\end{document}